\documentclass[11pt]{article}

\usepackage[utf8]{inputenc}
\usepackage[T1]{fontenc}
\usepackage{lmodern}
\usepackage[a4paper,margin=1in]{geometry}
\usepackage{xcolor}
\usepackage{graphicx}
\usepackage{booktabs}
\usepackage{array}
\usepackage{longtable}
\usepackage{enumitem}
\usepackage{listings}
\usepackage{fancyhdr}
\usepackage{titlesec}
\usepackage{tcolorbox}
\usepackage{amsmath}
\usepackage[hidelinks]{hyperref}
\usepackage{xurl}

% ---- palette (matches achernar_audit) ----
\definecolor{acnavy}{HTML}{1F2A44}
\definecolor{acblue}{HTML}{2D6CDF}
\definecolor{acteal}{HTML}{0E7C7B}
\definecolor{acgray}{HTML}{5B6472}
\definecolor{aclight}{HTML}{EEF2F8}
\definecolor{acred}{HTML}{C0392B}
\definecolor{acamber}{HTML}{B9770E}
\definecolor{acgreen}{HTML}{1E8449}
\definecolor{accode}{HTML}{F4F6FB}

\hypersetup{colorlinks=true,linkcolor=acblue,urlcolor=acblue}

\pagestyle{fancy}
\fancyhf{}
\setlength{\headheight}{14pt}
\addtolength{\topmargin}{-2pt}
\renewcommand{\headrulewidth}{0.4pt}
\lhead{\textcolor{acgray}{\small Achernar — Commercial Strategy \& Revenue Model}}
\rhead{\textcolor{acgray}{\small\thepage}}
\lfoot{\textcolor{acgray}{\small Confidential — Internal}}

\titleformat{\section}{\Large\bfseries\color{acnavy}}{\thesection}{0.6em}{}
\titleformat{\subsection}{\large\bfseries\color{acblue}}{\thesubsection}{0.6em}{}
\titleformat{\subsubsection}{\normalsize\bfseries\color{acteal}}{\thesubsubsection}{0.6em}{}

\newtcolorbox{keybox}[1][]{colback=aclight,colframe=acblue,boxrule=0.8pt,
  left=8pt,right=8pt,top=6pt,bottom=6pt,#1}
\newtcolorbox{warnbox}[1][]{colback=white,colframe=acamber,boxrule=0.8pt,
  left=8pt,right=8pt,top=6pt,bottom=6pt,#1}
\newtcolorbox{redbox}[1][]{colback=white,colframe=acred,boxrule=0.8pt,
  left=8pt,right=8pt,top=6pt,bottom=6pt,#1}

\lstset{
  basicstyle=\ttfamily\footnotesize,
  backgroundcolor=\color{accode},
  breaklines=true, frame=single, rulecolor=\color{aclight},
  framesep=6pt, xleftmargin=4pt, xrightmargin=4pt,
  columns=fullflexible, keepspaces=true, showstringspaces=false,
}

\newcommand{\dollar}[1]{\$#1}

\begin{document}

\begin{titlepage}
\centering
\vspace*{2.5cm}
{\Huge\bfseries\color{acnavy} Achernar\par}
\vspace{0.6cm}
{\LARGE\color{acblue} Commercial Strategy, Market Position\par}
{\LARGE\color{acblue} and Revenue Model\par}
\vspace{1.2cm}
{\large\color{acgray} Autonomy-First Launch Path\par}
\vspace{2.5cm}
{\large Ibrahim Shahid\par}
{\normalsize\color{acgray} Lead Architect\par}
\vspace{1.5cm}
{\normalsize\color{acgray} 7 August 2026\par}
\vfill
\begin{keybox}
\small
This document accompanies \texttt{achernar\_v4.pdf}. Market prices and unit
counts in Section~2 are drawn from public sources as of August 2026 and are
cited. All revenue figures in Sections~4--6 are \textbf{modelled outputs from
stated assumptions}, not forecasts. The assumptions --- not the totals --- are
the substance of this document; each is listed in Appendix~A so it can be
attacked individually.
\end{keybox}
\end{titlepage}

\tableofcontents
\newpage

% =====================================================================
\section{Executive Summary}

Achernar's commercial value does not lie in being a better capture application.
That market is served by mature free software and cannot be won on price. It
lies in \textbf{Section~7 and Section~8 of v4} --- learned closed-loop control
with calibrated uncertainty, bounded by a deterministic safety supervisor, and
validated through shadow mode against explicit promotion criteria, fed by a
single shared perception bus.

No competitor ships this. Every incumbent automates through deterministic
scripting and heuristics.

\begin{keybox}
\textbf{The thesis in one line:} do not sell imaging software; sell
\emph{unattended operation}, priced against the operator labour it removes.
\end{keybox}

\subsection{Headline numbers}

\begin{center}
\begin{tabular}{@{}lr@{}}
\toprule
\textbf{Quantity} & \textbf{Modelled value} \\
\midrule
Total addressable market (annual, all segments) & \dollar{8{,}490{,}000} \\
Year-5 ARR --- conservative & \dollar{126{,}700} \\
Year-5 ARR --- base & \dollar{708{,}800} \\
Year-5 ARR --- aggressive & \dollar{1{,}853{,}000} \\
Cumulative 5-year revenue --- base & \dollar{1{,}630{,}240} \\
Plausible acquisition value (base, 3--8$\times$ ARR) & \dollar{2.1}M--\dollar{5.7}M \\
Scopes needed to sustain a solo founder & 167 \\
OEM deals needed to sustain founder + 1 engineer & 1 \\
\bottomrule
\end{tabular}
\end{center}

\subsection{The honest shape of the opportunity}

\begin{warnbox}
This is a \textbf{small-market business}. The entire annual TAM across every
segment is under \dollar{9}M. Achernar is not a venture-scale company and should
not be planned as one. It is, on the base case, a sustainable one-to-three
person software business with a realistic \dollar{2}M--\dollar{6}M acquisition
outcome, and on the aggressive case a \dollar{5}M--\dollar{15}M one.
\end{warnbox}

The constraint is not execution and it is not the technology. It is that there
are only a few thousand robotic telescopes on Earth that anyone pays to operate.
Section~6 addresses the only lever that changes this: strategies that
\emph{expand} the market rather than dividing the existing one.

% =====================================================================
\section{Market Landscape}

\subsection{Price tiers (verified, August 2026)}

\begin{center}
\small
\begin{tabular}{@{}p{2.6cm}p{4.1cm}p{2.6cm}p{4.4cm}@{}}
\toprule
\textbf{Tier} & \textbf{Product} & \textbf{Price} & \textbf{Note} \\
\midrule
Free / OSS & NINA, Ekos/KStars, PHD2, Siril, ASTAP & \dollar{0} & Mature; broad hardware coverage \\
\addlinespace
Hardware bundle & ZWO ASIAIR Plus & \dollar{299} one-time & \textbf{ZWO gear only} \\
\addlinespace
Prosumer & Voyager & \texteuro129 + \texteuro29/yr & Markets ``AI driven automation''; DragScript / RoboTarget \\
\addlinespace
Prosumer & Sequence Generator Pro & Paid licence & Scripted sequencing \\
\addlinespace
\textbf{Institutional} & \textbf{ACP Expert} & \textbf{\dollar{1{,}495}} & Dispatch scheduler, PinPoint, remote observatory \\
\addlinespace
Institutional & TheSkyX Professional & Paid licence & TPoint, dome and camera control \\
\addlinespace
Service & iTelescope membership & \dollar{40}--\dollar{160}/mo & End-user access to hosted rigs \\
\addlinespace
Service & DSP Remote hosting & from \dollar{300}/mo & Per-rig hosting fee \\
\bottomrule
\end{tabular}
\end{center}

\vspace{0.4em}
Smart telescope hardware market: \textbf{\dollar{221}M (2025) $\rightarrow$
\dollar{360}M (2032)}, 7.28\% CAGR. The ZWO Seestar S50 alone sold
\textbf{over 50,000 units} before its early-2026 discontinuation, replaced by
the S50 Pro at \dollar{599}--\dollar{799}.

\subsection{The white space}

Every product in the table above automates through \textbf{scripts and rules}.
Voyager's DragScript is genuinely sophisticated, but it is deterministic
scripting with heuristic triggers, not learned control. Nothing on the market
ships:

\begin{itemize}[leftmargin=1.4em,itemsep=2pt]
  \item learned closed-loop control with \textbf{calibrated uncertainty};
  \item a \textbf{deterministic safety supervisor} bounding every model action;
  \item \textbf{shadow-mode evaluation} with explicit promotion criteria;
  \item \textbf{cross-subsystem shared perception} --- one vision loop feeding
        focus, filter wheel and mount from a single observation.
\end{itemize}

v4 §7 and §8 specify all four. This is a real gap in the market, not a
positioning exercise.

\subsection{Two structural openings}

\subsubsection*{Opening 1 --- ACP Expert at \dollar{1{,}495} is the price anchor}

ACP Expert sells at roughly ten times the prosumer tier and it sells on exactly
one property: unattended operation. Institutions demonstrably pay that figure
for automation built on two-decade-old dispatch scheduling. This validates the
pricing thesis for an autonomy layer far more strongly than any hobbyist
subscription analysis could.

\subsubsection*{Opening 2 --- ASIAIR is ZWO-only}

Achernar's INDI / ASCOM / direct-USB abstraction is vendor-neutral, which is the
one property ZWO structurally cannot copy without cannibalising camera sales.

\begin{redbox}
\textbf{Caveat.} \texttt{hardware/indi\_client.rs} is currently a 12-line stub.
The moat is specified and not yet built. Until it is, the vendor-neutrality
argument is a claim, not a capability.
\end{redbox}

% =====================================================================
\section{Repository State}

Verified against code on 7 August 2026, not inferred from documentation.

\subsection{Built --- 19,626 LOC Rust}

\begin{center}
\small
\begin{tabular}{@{}lrp{8.4cm}@{}}
\toprule
\textbf{Module} & \textbf{LOC} & \textbf{State} \\
\midrule
\texttt{hardware/} & 3{,}547 & camera, focuser, mount, filter wheel, OLS FFI, gphoto2 DSLR, raw decode, discovery \\
\texttt{vision/} & 3{,}115 & \textbf{vision loop at 1--5\,Hz, SharedObservation bus (§8.1 complete)}, classical HFR / star detection, background, trailing, FQE runner \\
\texttt{api/} & 3{,}021 & camera, focus, focuser, mount, websocket \\
\texttt{focus/} & 2{,}553 & agent state machine, classical solver, backlash, drift, features, policy \\
\texttt{storage/} & 1{,}944 & 15 repositories, migrations 001--005 \\
\texttt{mount/} & 693 & guiding bandit, dither policy, plate-solve trigger, pointing model, safety \\
\texttt{filter\_wheel/} & 466 & SNR scorer, bandit, planner, controller, safety \\
\texttt{control/} & 436 & \textbf{safety guard, action arbiter, control supervisor, ownership} \\
\texttt{models/} & 374 & registry with SHA-256 and \texttt{runtime\_eligible} gating \\
\bottomrule
\end{tabular}
\end{center}

\vspace{0.5em}
Plus \texttt{ml/} at 4,501 LOC (excluding virtualenv): dataset builders, FQE and
offset training, ONNX export, evaluation, replay harness, synthetic PSF,
AstroM3 loader with checkpoints downloaded.

\begin{keybox}
The v4 §7/§8 control scaffolding is \textbf{essentially complete}. This is the
commercially differentiating layer, and it exists.
\end{keybox}

\subsection{The single blocker: no trained models, no data}

\begin{redbox}
\begin{itemize}[leftmargin=1.2em,itemsep=2pt]
  \item \texttt{models/focus\_quality/1.0.0/metadata.json} declares
        \texttt{"trained": false}, \texttt{"weights": "random"},
        \texttt{"runtime\_eligible": false},
        \texttt{"promotion\_status": "never --- integration placeholder only"}.
  \item \texttt{focus\_offset/}, \texttt{cloud\_classifier/} and
        \texttt{trailing\_detector/} carry \textbf{no \texttt{training\_info}
        block at all}, so they fail \texttt{is\_runtime\_eligible()}
        (\texttt{models/model\_registry.rs:55}) as well.
  \item \texttt{frames} table: \textbf{0 rows}. \texttt{focus\_samples}: 7.
        \texttt{ml/data/}: a README with four unchecked TODOs.
\end{itemize}
\end{redbox}

The architecture is correctly refusing to run itself. Both
\texttt{model\_registry.rs:55} and \texttt{vision/fqe\_runner.rs:325} gate on
\texttt{runtime\_eligible}, so every model path stays dark until real weights
exist. That is the design working as intended --- but it means \textbf{the data
flywheel has never turned once}.

\subsection{Missing, on the critical path}

\begin{itemize}[leftmargin=1.4em,itemsep=2pt]
  \item \texttt{processing/} is \textbf{58 LOC total}. \texttt{debayer.rs},
        \texttt{calibration.rs} and \texttt{stacking.rs} are four-line
        \texttt{pub struct} declarations with TODOs.
        \texttt{plate\_solver.rs} always returns \texttt{Unsupported}.
  \item \texttt{hardware/indi\_client.rs} --- 12-line stub (the moat, above).
  \item \texttt{ai/} (18 LOC) is dead code, superseded by \texttt{vision/}.
        Delete it.
\end{itemize}

% =====================================================================
\section{Revenue Model}

\subsection{Method}

Bottom-up by segment: universe size $\times$ annual revenue per unit $\times$
penetration. Universe sizes are grounded in the public data of Section~2 where
it exists and estimated where it does not; every estimate is flagged in
Appendix~A. Penetration rates are judgement.

\[
\text{ARR}_{\text{segment}} = N_{\text{universe}} \times p_{\text{penetration}} \times \text{ARPU}_{\text{annual}}
\]

\subsection{Segments and total addressable market}

\begin{center}
\small
\begin{tabular}{@{}lrrr@{}}
\toprule
\textbf{Segment} & \textbf{Universe} & \textbf{ARPU/yr} & \textbf{TAM/yr} \\
\midrule
Prosumer autonomy licence & 20{,}000 & \dollar{110} & \dollar{2{,}200{,}000} \\
Hosting operator (per scope) & 3{,}500 & \dollar{540} & \dollar{1{,}890{,}000} \\
Institutional site licence & 800 & \dollar{3{,}000} & \dollar{2{,}400{,}000} \\
OEM autonomy stack & 8 & \dollar{250{,}000} & \dollar{2{,}000{,}000} \\
\midrule
\textbf{Total} & & & \textbf{\dollar{8{,}490{,}000}} \\
\bottomrule
\end{tabular}
\end{center}

\vspace{0.3em}
\noindent\textbf{ARPU derivations.} Prosumer: \dollar{149} perpetual amortised
plus \dollar{39}/yr renewal, priced just above Voyager. Hosting:
\dollar{45}/scope/month. Institutional: anchored on ACP Expert's \dollar{1{,}495}
perpetual, converted to a \dollar{3{,}000}/yr site licence covering multiple
instruments. OEM: flat licence plus royalty.

\subsection{Year-5 steady-state scenarios}

\begin{center}
\small
\begin{tabular}{@{}lrrr@{}}
\toprule
\textbf{Segment} & \textbf{Conservative} & \textbf{Base} & \textbf{Aggressive} \\
\midrule
Prosumer autonomy licence & \dollar{22{,}000} & \dollar{88{,}000} & \dollar{176{,}000} \\
Hosting operator (per scope) & \dollar{56{,}700} & \dollar{226{,}800} & \dollar{567{,}000} \\
Institutional site licence & \dollar{48{,}000} & \dollar{144{,}000} & \dollar{360{,}000} \\
OEM autonomy stack & \dollar{0} & \dollar{250{,}000} & \dollar{750{,}000} \\
\midrule
\textbf{Total ARR} & \textbf{\dollar{126{,}700}} & \textbf{\dollar{708{,}800}} & \textbf{\dollar{1{,}853{,}000}} \\
Customers / units & 321 & 1{,}269 & 2{,}773 \\
\midrule
\multicolumn{4}{@{}l@{}}{\footnotesize Penetration: prosumer 1/4/8\%; hosting 3/12/30\%; institutional 2/6/15\%; OEM 0/1/3 deals.} \\
\bottomrule
\end{tabular}
\end{center}

\subsection{Five-year ramp (base case)}

\begin{center}
\small
\begin{tabular}{@{}lrrp{7.2cm}@{}}
\toprule
\textbf{Year} & \textbf{ARR} & \textbf{Cumulative} & \textbf{Milestone} \\
\midrule
Y1 & \dollar{21{,}264} & \dollar{21{,}264} & Shadow-mode validation; design partners; no revenue at scale \\
Y2 & \dollar{106{,}320} & \dollar{127{,}584} & First paid hosting sites; prosumer beta \\
Y3 & \dollar{283{,}520} & \dollar{411{,}104} & Institutional entry; prosumer GA \\
Y4 & \dollar{510{,}336} & \dollar{921{,}440} & OEM deal lands; hosting expands \\
Y5 & \dollar{708{,}800} & \dollar{1{,}630{,}240} & Steady state \\
\bottomrule
\end{tabular}
\end{center}

\subsection{Breakeven thresholds}

The most useful table in this document, because it converts strategy into a
countable target.

\begin{center}
\small
\begin{tabular}{@{}lrrrrr@{}}
\toprule
\textbf{Cost base} & \textbf{Burn/yr} & \textbf{Scopes} & \textbf{Sites} & \textbf{Prosumers} & \textbf{OEM deals} \\
\midrule
Solo founder, full-time & \dollar{90{,}000} & 167 & 30 & 818 & 0.4 \\
Founder + 1 engineer & \dollar{240{,}000} & 444 & 80 & 2{,}182 & 1.0 \\
Team of 4 & \dollar{620{,}000} & 1{,}148 & 207 & 5{,}636 & 2.5 \\
\bottomrule
\end{tabular}
\end{center}

\begin{keybox}
\textbf{167 scopes.} That is the entire requirement to make Achernar a
full-time job at \dollar{45}/scope/month. Four mid-size hosting sites. This is
the single most important number in the document --- it is small enough to be
reachable and it does not depend on the prosumer market at all.
\end{keybox}

% =====================================================================
\section{Maximising Revenue: Strategies Ranked}

Ranked by expected value, which is not the same as ranked by ceiling.

\subsection{1. Per-scope recurring revenue from hosting operators}

\textbf{Best revenue per unit of effort in the near term.}

Commercial hosting sites charge from \dollar{300}/month per rig. At
\dollar{45}/scope/month you are asking for 15\% of their gross on a line item
that removes operator labour --- an easy internal justification. Sales cycle is
short because the buyer is an owner-operator, not a procurement committee.

\begin{itemize}[leftmargin=1.4em,itemsep=2pt]
  \item \textbf{Why it wins:} recurring, scales with fleet size, and the
        customer's own economics do the selling.
  \item \textbf{Target:} 30--60 commercial facilities worldwide. Perhaps 20
        conversations to first contract.
  \item \textbf{Requirement:} INDI client (they run heterogeneous gear), and
        demonstrated fallback and clamp rates in shadow mode.
\end{itemize}

\subsection{2. OEM licensing of the autonomy stack}

\textbf{Highest ceiling per unit of effort.}

Every non-ZWO vendor --- Player One, QHYCCD, ToupTek, Pegasus Astro,
PrimaLuceLab, Celestron, Sky-Watcher --- needs software parity against ASIAIR
and none of them want to build a nine-model control stack. You sell ONNX
bundles, the safety supervisor and an SDK; they sell hardware.

\textbf{One deal at \dollar{250{,}000}/yr sustains a founder and an engineer
outright.} That is the highest-leverage single transaction available.

\subsection{3. Smart-telescope attach --- the only path with unit scale}

The Seestar S50 sold \textbf{over 50,000 units}. That is the only genuinely
large number anywhere in this analysis. Powering a credible number-two vendor:

\begin{center}
\small
\begin{tabular}{@{}lrr@{}}
\toprule
\textbf{Model} & \textbf{15,000 units} & \textbf{30,000 units} \\
\midrule
One-time royalty @ \dollar{15}/unit & \dollar{225{,}000} & \dollar{450{,}000} \\
Subscription attach @ \dollar{5}/mo, 12\% take & \dollar{108{,}000}/yr & \dollar{216{,}000}/yr \\
Subscription attach @ \dollar{5}/mo, 25\% take & \dollar{225{,}000}/yr & \dollar{450{,}000}/yr \\
\bottomrule
\end{tabular}
\end{center}

\subsection{4. Vertical integration --- operate, do not licence}

\begin{center}
\small
\begin{tabular}{@{}lr@{}}
\toprule
Hosting revenue per rig per month & \dollar{300} \\
Rigs at one mid-size site & 40 \\
Site gross revenue per year & \dollar{144{,}000} \\
Same site licensed at \dollar{45}/scope/month & \dollar{21{,}600}/yr \\
\midrule
\textbf{Revenue multiple of owning vs. licensing} & \textbf{6.7$\times$} \\
\bottomrule
\end{tabular}
\end{center}

\vspace{0.3em}
Owning the site captures 6.7$\times$ the revenue of licensing software to it,
and the autonomy stack is precisely what lets one operator run a larger fleet.
The cost is capital, land, weather risk and operations --- a fundamentally
different company. Correct as a year-4+ move, wrong as a starting point.

\subsection{5. Institutional site licences}

ACP-anchored at \dollar{3{,}000}/yr. Slow (6--18 month cycles, grant budgets)
but extremely sticky and reference-generating. 30 sites sustains a solo founder.
Pursue in parallel with hosting; do not make it the first market, because the
cycle length will starve you.

\subsection{6. Prosumer --- deliberately not a revenue line}

\begin{keybox}
\textbf{Recommendation: make the prosumer tier free.}

At \dollar{88{,}000} base-case ARR it is the smallest segment in the model, it
carries the highest support burden per dollar, and it competes directly with
free incumbents. Its actual value is as \textbf{the data flywheel} --- the
diverse rigs, skies and failure modes that every model in §7.6 needs --- and as
the credibility that makes the OEM and institutional conversations possible.

Charging \dollar{110}/yr for it trades a \dollar{2}M strategic asset for a
\dollar{88}k revenue line.
\end{keybox}

% =====================================================================
\section{The Maximum-Money Thesis}

\subsection{The binding constraint is market size, not share}

Every strategy in Section~5 divides a TAM under \dollar{9}M. Even perfect
execution across all four segments yields under \dollar{2}M ARR. Optimising
pricing or penetration inside that box cannot produce a large outcome.

\begin{keybox}
The only strategies that materially change the ceiling are those that
\textbf{expand the number of telescopes operating unattended}, rather than
capturing a larger share of the few thousand that already do.
\end{keybox}

\subsection{Three market-expanding moves}

\subsubsection*{(a) Drive the cost of unattended operation below the hosting price}

Remote hosting exists because operating a remote rig requires a human. If
Achernar's autonomy stack makes a \dollar{2{,}000} backyard rig genuinely
unattended --- roll-off roof, weather abort, automatic recovery, no laptop, no
babysitting --- then the addressable population is not 3,500 hosted scopes, it
is every serious imager with a permanent pier. That is a two-order-of-magnitude
expansion of the same product.

\subsubsection*{(b) Own the smart-telescope autonomy layer}

50,000 Seestar units in roughly two years is the only real volume in this
market, and that volume exists because the product removed operation from the
user entirely. The second generation of smart telescopes will compete on
capability under autonomy --- longer integrations, mosaics, multi-night
projects, transient response. That is exactly the §7 stack. Being the layer
inside three vendors' products is worth more than being an application.

\subsubsection*{(c) Sell the fleet layer, then become the fleet}

Land hosting operators as customers first (Section~5.1), learn the operational
economics from inside, then vertically integrate at 6.7$\times$ the revenue per
site. The software business is the reconnaissance for the operating business.

\subsection{The realistic liquidity event}

Niche vertical software in hardware-adjacent markets clears at roughly
3--6$\times$ ARR; a strategic buyer defending a platform position pays more.

\begin{center}
\small
\begin{tabular}{@{}lrrr@{}}
\toprule
\textbf{ARR at exit} & \textbf{3$\times$} & \textbf{5$\times$} & \textbf{8$\times$ (strategic)} \\
\midrule
\dollar{250{,}000} & \dollar{750{,}000} & \dollar{1{,}250{,}000} & \dollar{2{,}000{,}000} \\
\dollar{709{,}000} (base) & \dollar{2{,}127{,}000} & \dollar{3{,}545{,}000} & \dollar{5{,}672{,}000} \\
\dollar{1{,}853{,}000} (aggressive) & \dollar{5{,}559{,}000} & \dollar{9{,}265{,}000} & \dollar{14{,}824{,}000} \\
\dollar{3{,}000{,}000} & \dollar{9{,}000{,}000} & \dollar{15{,}000{,}000} & \dollar{24{,}000{,}000} \\
\bottomrule
\end{tabular}
\end{center}

\vspace{0.3em}
Plausible acquirers: ZWO (defending ASIAIR against a vendor-neutral rival),
Celestron or Sky-Watcher (needing a Seestar answer), Software Bisque or
Diffraction Limited (modernising an ageing institutional stack), or a hosting
operator buying its own supply chain.

\begin{warnbox}
\textbf{Accumulating ARR is probably not how this returns the most money.} A
strategic acquisition at 5--8$\times$, driven by a buyer who needs the autonomy
stack to defend a hardware position, is both faster and larger than organic
growth inside a \dollar{9}M TAM. Build accordingly: prioritise the
\emph{demonstrable, differentiated capability} and the reference customers that
prove it, over revenue maximisation per seat.
\end{warnbox}

% =====================================================================
\section{What Kills Each Path}

\begin{center}
\small
\begin{tabular}{@{}p{3.6cm}p{5.2cm}p{5.2cm}@{}}
\toprule
\textbf{Path} & \textbf{Primary risk} & \textbf{Mitigation} \\
\midrule
Hosting operators & One mount crash destroys the reference account and the reputation with it & Safety supervisor is already built; never ship above supervised-auto without shadow-mode promotion evidence \\
\addlinespace
OEM licensing & ZWO bundles a worse autonomy layer for free & Vendor neutrality --- build the INDI client; ZWO cannot follow without cannibalising camera sales \\
\addlinespace
Institutional & 6--18 month cycles starve a solo operation & Never make it the first market; run in parallel with hosting revenue \\
\addlinespace
Prosumer & Competes with free; support burden exceeds revenue & Make it free by design; harvest data instead \\
\addlinespace
Vertical integration & Capital, weather, land, operations --- a different company & Year-4+ move only, funded from software revenue \\
\addlinespace
All paths & MIT licence on a system that moves hardware & Relicense the control and ML layers before any commercial release \\
\bottomrule
\end{tabular}
\end{center}

\begin{redbox}
\textbf{Licensing.} The README declares MIT. A permissive licence on a
closed-loop control system that physically drives a mount is both a commercial
error (anyone may fork the moat) and a liability exposure. Resolve before any
external distribution.
\end{redbox}

% =====================================================================
\section{Implementation Sequence}

Commercial strategy is downstream of one fact from Section~3: \textbf{the
\texttt{frames} table has zero rows}. Nothing in this document is reachable
until the flywheel turns once.

\begin{enumerate}[leftmargin=1.6em,itemsep=4pt]
  \item \textbf{Compile.} \texttt{HANDOFF-next-session.md} records that the
        capture rework was never compiled --- no Rust toolchain in the sandbox.
        Run \texttt{cargo build \&\& cargo test} locally. Expect signature
        mismatches. Nothing below is real until this passes.

  \item \textbf{Turn the flywheel once.} Capture end-to-end until
        \texttt{frames} receives its first row.

  \item \textbf{Debayer and calibration} (\texttt{processing/}, currently 8 LOC
        of TODO). FQE consumes $64\times64$ crops; without debayering, an OSC or
        DSLR frame yields HFR measured across a Bayer mosaic --- corrupt labels.
        This precedes any training.

  \item \textbf{Focus sweep harness $\rightarrow$ first dataset.} Automated ramp
        across focuser positions, capture at each, classical HFR as label
        (§9.1 bootstrapping). One clear night yields a few thousand labelled
        crops. \texttt{ml/data\_pipeline/build\_focus\_quality\_dataset.py}
        already exists to consume it.

  \item \textbf{Train FQE for real.} The chain already exists end to end:

\begin{lstlisting}
ml/train/train_focus_quality.py
  -> ml/train/calibrate_uncertainty.py
  -> ml/export/export_quality_onnx.py
  -> ml/export/write_model_metadata.py
\end{lstlisting}

        Then flip metadata to \texttt{trained: true},
        \texttt{runtime\_eligible: true}, \texttt{promotion\_status:
        "shadow"}. That single change lights up the shadow path already
        wired into the vision loop at line~120.

  \item \textbf{Shadow-mode evaluation.} \texttt{ml/eval/replay\_offset\_policy.py}
        exists. Compare FQE against classical HFR over logged sessions; apply
        §7.10 promotion criteria.

  \item \textbf{INDI client.} The vendor-neutrality moat, and the prerequisite
        for every hosting-operator conversation in Section~5.1.
\end{enumerate}

\begin{keybox}
\textbf{Why FQE first.} Not because focus matters most, but because it is the
\emph{only} model where classical bootstrapping (§9.1) produces labels for free.
Every other model in §7.6 requires human labelling or nights of policy rollout.
Close one full lap --- train, export, shadow, promote --- and every remaining
model is the same lap with a different label source. Prove the lap once,
carefully.
\end{keybox}

\subsection{Commercial milestones mapped to technical ones}

\begin{center}
\small
\begin{tabular}{@{}p{5.4cm}p{8.6cm}@{}}
\toprule
\textbf{Technical milestone} & \textbf{Commercial unlock} \\
\midrule
First FQE promotion out of shadow mode & Evidence that the promotion machinery works --- the core claim of the pitch \\
Fallback and clamp rates logged over 20+ nights & The only data a hosting operator will actually ask for \\
INDI client operational & Hosting and institutional conversations become possible at all \\
Three models promoted, two rigs & First paid design partner \\
Unattended multi-night session, zero interventions & The demo that closes OEM conversations \\
\bottomrule
\end{tabular}
\end{center}

% =====================================================================
\newpage
\appendix
\section{Assumptions Register}

Each row is a load-bearing assumption. Attack these rather than the totals.

\begin{center}
\small
\begin{longtable}{@{}p{4.4cm}rp{7.6cm}@{}}
\toprule
\textbf{Assumption} & \textbf{Value} & \textbf{Basis / confidence} \\
\midrule
\endhead
Paid-automation user base & 20{,}000 & \textcolor{acamber}{Estimated.} Voyager + SGPro class installed base. No public figure; could be 10k or 40k. \\
Commercially hosted rigs worldwide & 3{,}500 & \textcolor{acamber}{Estimated} from 30--60 known facilities at 20--100 rigs each. \\
Small-aperture robotic + teaching observatories & 800 & \textcolor{acamber}{Estimated} from 350+ US and 600+ non-US professional observatories, small-aperture subset. \\
Non-ZWO OEM candidates & 8 & \textcolor{acgreen}{High confidence.} Enumerable vendor list. \\
Hosting ARPU & \dollar{45}/scope/mo & \textcolor{acamber}{Judgement.} 15\% of the \dollar{300}/mo hosting fee. Untested. \\
Institutional ARPU & \dollar{3{,}000}/yr & \textcolor{acgreen}{Anchored} on verified ACP Expert \dollar{1{,}495} perpetual. \\
OEM deal value & \dollar{250{,}000}/yr & \textcolor{acred}{Low confidence.} No comparables in this market. \\
Hosting fee per rig & \dollar{300}/mo & \textcolor{acgreen}{Verified} --- DSP Remote published pricing. \\
Seestar S50 units & 50{,}000+ & \textcolor{acgreen}{Verified} from multiple retail and review sources. \\
Smart telescope market & \dollar{221}M $\rightarrow$ \dollar{360}M & \textcolor{acgreen}{Verified} --- published market research, 7.28\% CAGR. \\
Penetration rates & 1--30\% & \textcolor{acred}{Judgement only.} The weakest input in the model. \\
Exit multiple & 3--8$\times$ ARR & \textcolor{acamber}{Sector-typical} for niche vertical software. \\
\bottomrule
\end{longtable}
\end{center}

\begin{warnbox}
\textbf{Sensitivity.} Base-case ARR is dominated by two inputs: the OEM deal
(\dollar{250}k of \dollar{709}k, 35\%) and hosting penetration. If no OEM deal
lands, base case falls to \dollar{459}k. If hosting penetration is 5\% rather
than 12\%, it falls a further \dollar{132}k. The model is therefore best read as:
\textbf{\dollar{300}k--\dollar{700}k is the realistic band}, with everything above
it contingent on a single OEM relationship.
\end{warnbox}

\section{Sources}

\begin{itemize}[leftmargin=1.4em,itemsep=1pt]
\small
  \item ACP Expert pricing --- \url{https://acpx.dc3.com/} and OPT Telescopes
  \item Voyager pricing --- \url{https://software.starkeeper.it/voyager-advanced/}
  \item DSP Remote hosting --- \url{https://www.dspremote.com/plans-and-pricing/}
  \item iTelescope membership --- \url{https://www.itelescope.net/}
  \item Smart telescope market size --- \url{https://www.verifiedmarketreports.com/product/smart-telescope-market/}
  \item Seestar S50 units and discontinuation --- \url{https://www.telescopeadvisor.com/seestar-s50-discontinued-what-to-buy-instead/}
  \item Observatory counts --- \url{https://www.go-astronomy.com/observatories.htm}
  \item Las Cumbres network --- \url{https://lco.global/about/}
  \item NINA --- \url{https://nighttime-imaging.eu/}
\end{itemize}

\end{document}
